\summarychapter{40}{A Bridge for Intelligent Moderns}{Protestant Liberalism}

\begin{inbrief}
Protestant liberalism tried to lead the churches into the world of modern science, philosophy and
history, so that a person could be, in Harry Emerson Fosdick's words, ``both an intelligent modern
and a serious Christian.'' The chapter pictures it as a bridge resting on two towers, modern thought
(evolution and biblical criticism) and Christian experience (Schleiermacher and Ritschl), and shows
the change at work among New England Congregationalists.
\question{How can the faith be restated for a new intellectual world without being bent out of shape or lost?}
\end{inbrief}

On his eighty-fifth birthday in 1920 Lyman Abbott, who in the 1890s had been among the most influential preachers in America, looked back on his Puritan childhood, when he thought of God as ``a kind of awful and
omnipresent police justice'' and of himself as ``a scared culprit'' who did not clearly know why he
deserved punishment. He had long since given up both pictures. Abbott was typical of many ministers
raised in pious Protestant homes who studied in Germany, or in American seminaries that prized German
scholarship, and adopted liberal convictions. The twentieth century has been unkind to the liberal
creed, yet every major Protestant denomination still shows its mark, and Sydney Ahlstrom judged that
the liberals set off the most fundamental controversy in the churches since the Reformation.

\ssection{The Aims of Protestant Liberalism}
The problem is as old as Christianity. Paul tried to restate the faith for a new world of thought and
succeeded; the Gnostics tried and failed. On liberalism no final verdict has been given, though Christian
opinion leans heavily toward failure. H.~Richard Niebuhr put its irony best: ``a God without wrath
brought men without sin into a kingdom without judgment through the ministrations of a Christ without
a Cross.'' Many deny that it is a theology at all and call it a spirit; Henry Sloane Coffin of Union
Seminary called it the spirit that puts truth above everything and so demands liberty to debate,
publish and pursue it.

\begin{sidebox}{Theological Liberalism}
Theological liberalism takes a progressive ethic from the Christian message and rejects the doctrines in which Christians have traditionally stated that message. Its perfect example is Adolf von
Harnack (1851--1930), a pioneer in New Testament studies, theology, ethics and historical theology,
which he revolutionized. He argued that Jesus was misunderstood and misrepresented in most Christian
sources, the Gospels included. Jesus was a simple Hebrew prophet calling people to moral awakening,
and his real message can be recovered by stripping off the layers of tradition, the husk from the
kernel. Harnack's Jesus proclaimed a kingdom that is moral and present, not a supernatural upheaval still to come; his themes were the fatherhood of God, the brotherhood of mankind and
the infinite value of every soul, and called his followers to a higher righteousness rooted in God's
command to love.
\end{sidebox}

Shelley holds that the liberal spirit does lead to identifiable convictions and that spirit and
convictions together make up liberalism. He likens it to a suspension bridge. One tower stands on
modern thought, the other on Christian experience, and the soil shifts under both. Those who use the
bridge disagree about which side is safer, which is why Kenneth Cauthen finds two types. Evangelical
liberals were serious Christians looking for a theology that could serve intelligent moderns, and
they trusted the tower of experience. Modernistic liberals were intelligent moderns who hoped to be
counted Christians in some sense, and they trusted the tower of modern thought.

Liberals believed that theology must come to terms with modern science if it hoped to hold
intelligent people. They refused to accept beliefs on authority alone. Faith must pass the tests of
reason and experience; the human mind can think God's thoughts after him, and the best clues to God's
nature are intuition and reason. A Christian should be open to truth from any source, and an
unexamined faith is not worth having. In yielding so fully to the modern mind they accepted its
assumption that the universe is one harmonious machine or growing organism. The biblical account
recognizes distinct orders---matter, plants, animals, man and God---but liberal theology pressed
for continuity. The lines it blurred were many: revealed and natural religion, Christianity and the other faiths, the saved and the lost, Christ and the rest of humanity, even God and man.

Two terms are crucial. Immanence is God's presence in the world, working through nature; its extreme
is pantheism, which identifies God and the world. Transcendence is God's reality apart from the
world; its extreme is deism, with God as separate as a watchmaker from his watch. Liberals thought
the idea of a God somewhere beyond the universe unacceptable to modern people (a view of God, Shelley
notes, that owes more to deism than to real Christian theology). They linked the divine with human
consciousness and called the life running through nature and man God. This fitted science. God had
not broken through the clouds to create the world but had worked for ages through natural law. As
the poet said, ``Some call it evolution, and others call it God.'' Evolution holds that natural selection has produced every complex organism from simpler ones, so that no species
is fixed. In 1785 James Hutton attributed the earth's development to natural causes, and Sir Charles
Lyell's \emph{Principles of Geology} (1830) confirmed that the earth's surface was formed by natural
causes over a vast time, the time that any theory of evolution needs.

\ssection{Enter Charles Darwin}
Charles Darwin (1809--1882) studied medicine, prepared at Cambridge for the ministry and became a
naturalist. From 1831 to 1836 he worked on the specimens gathered during the surveying voyage of the
\emph{Beagle} along the coast of South America. \emph{The Origin of Species} (1859) contended that
species have been modified over a long course of descent, chiefly through the natural selection of
many slight, favorable variations. It revolutionized ideas about the origin of life. \emph{The
Descent of Man} (1871) applied natural selection to human beings and concluded that their ancestors
were probably monkey-like animals. Many religious people were thrown on the defensive. Some rejected
the new views outright: if man is not specially created and fallen, what need is there of Christ's
salvation? Others tried to reconcile faith and evolution, and liberals came to think that evolution
supplements Christianity, growth and development being God's way of revealing himself. In 1892 Lyman
Abbott, minister of Plymouth Church in Brooklyn, wrote \emph{The Evolution of Christianity} and spoke
of the evolution of the Bible, the church and even the soul.

Serious as it was, the challenge of science was secondary to that of history. Science could question
God's rule over the physical world; historical criticism went straight to God's revelation in the
Bible. ``Criticism'' names a method, the subjecting of material to careful scrutiny, for example
asking whether there is historical evidence for the church's teaching that the Son is God or whether
it rests on tradition. The movement assumed that it was itself free of corrupting traditions, an
assumption harder to believe today. Lower criticism tries to reconstruct the exact wording of the
lost originals from the many surviving manuscripts, and conservatives, who care greatly about the
exact wording, have contributed much to it. Higher criticism is another matter. It asks about meaning
and whether the events reported really happened, and studies the background, authorship and setting
of each text. The method was not new, but in the nineteenth century it was applied to the Bible as
to any other ancient book. In practice, Shelley objects, the critics did not treat it like any other
book: many assumed that God cannot act in history, and others that the Bible's story is untrue unless
confirmed by an independent source.

The conclusions shook orthodoxy. The books of Moses were said to be the product of centuries, four
great literary traditions combined in the course of Israel's history; some prophecies seemed true
because they were written after the events; the accuracy of the Gospels, especially John, was
challenged. A central concern was the search for the ``historical Jesus,'' on the assumption that
Jesus as he lived differed from the Jesus of the Gospels and that authentic sayings and deeds must be
sifted from later additions by the early church. David Friedrich Strauss (1809--1874) rejected the
earlier approaches: the straightforward supernaturalist reading, the face-saving rationalism that
found natural explanations for miracles (Jesus walking on a jetty, not on the water), and the notion
of deliberate fraud by the evangelists. The Gospels, he argued, reflect the mythological outlook of
the ancients. Myth was the vehicle by which ancient people grasped great persons and events, and a
modern reader who knows this can assess the ideas the myths carry.

\ssection{The Impact of Biblical Criticism}
More important than any single conclusion was that criticism put the general trustworthiness of the
Bible in question. Liberals welcomed it. They no longer had to defend the whole Bible as the
infallible Word of God, or a God who ordered Israel to kill its enemies to the last woman and child
or sent bears to maul children who mocked a prophet. Criticism showed, they said, that revelation itself had evolved. Its first notions were crude and bloodthirsty; it advanced as
the Jews slowly grasped a righteous God served by doing justly, loving mercy and walking humbly, and
reached fulfillment in Jesus, who shows God as the loving Father of all.

When the traditional doctrines no longer gave assurance, liberals turned to the other tower,
Christian experience. Early in the century the movement called Romanticism, usually understood as a
protest against the rationalism of the Enlightenment, discerned a dynamic energy surging in nature,
not irrational but beyond reason; Romantic painters showed it as a great storm, dangerous and
fascinating at once. It is met by intuition, not logical deduction, and to meet it is of first
importance. Why trouble about external creeds when so intimate a certainty rules every soul?
``Closer is He than breathing,'' wrote Tennyson, ``and nearer than hands and feet.'' In 1891 Lewis
F.~Stearns of Bangor Seminary told the International Congregational Council in London that criticism
and rationalism could no more disturb the Christian realities in the believer's experience than the
surf on the coast of Maine could overthrow its granite cliffs.

\ssection{Influential Spokesmen for Liberalism}
The two chief spokesmen for experience were German theologians, Friedrich Schleiermacher
(1768--1834) of the new University of Berlin and Albrecht Ritschl (1822--1889) of Bonn and
G\"ottingen. Following William Hordern's account, Schleiermacher tried to rehabilitate religion among
intellectuals. Debates over proofs of God and abstract doctrine are at best secondary. The heart of
religion is the awareness of our absolute dependence before the grandeur of God. Cultivated people
were right to reject the abstractions of Protestant Scholasticism; their Romantic sensibilities were
already religious, and Christianity rightly understood explains that experience better than any
other religion. After him it became common to regard the religions as differing expressions of one
inner experience shared by all people. Christ's uniqueness lies not in a doctrine or a miraculous
origin such as the virgin birth: ``The real miracle is Jesus himself,'' the man with supreme
God-consciousness. We glimpse God now and then, while he knew him fully; our obedience comes and goes, while his was entire. Having that knowledge he can communicate God-consciousness to others, and the
church is the living witness that through contact with his life people have come to a vital
relationship with God and to brotherly living. Schleiermacher is ``the father of modern theology''
because he moved the basis of Christian faith from the Bible to religious experience.

Ritschl, the most influential theologian of the late century and the chief teacher of American
liberals, fixed that experience on the historical Jesus. Religion must be practical. To be saved is
not to go to heaven at death, a theoretical matter, but to live a new life, free of sin, selfishness,
fear and guilt. Christianity must be built on fact, so he welcomed the search for the historical
Jesus; the great fact is the impact Jesus has made on the church through the centuries. Nature speaks
ambiguously of God. We find him in history, where movements arise devoted to the values that give
life meaning. Religion rests on values, not on scientific truth: science gives the facts, and
religion weighs them. Man, though a product of nature and evolution, has a sense of values, and God
is the postulate needed to explain it. Many found this helpful, for it seemed to free faith from the
damage done by science and history. Criticism could settle authorship, date and meaning as matters of
fact, but it cannot rule on values. If it denies the miracles, the virgin birth and the preexistence
of Jesus, he is no less valuable. His divinity rests on his being the source of a value-creating
movement. He embodied such high ethical ideals that we are inspired to live as he did, and he is
divine in the sense that he does for us what God does, making us conscious of the highest in life.
From him comes the church, a value-creating community and the spearhead of a society inspired by
love, the kingdom of God on earth.

\ssection{An Example of Liberal Change}
Liberalism challenged orthodox bodies all over Europe and North America. The Congregationalists of
the United States serve as an example. Among them it appeared as the New Theology, and its advocates
came from within New England Calvinism: Theodore Thornton Munger and Newman Smyth, ministers in New
Haven; George Angier Gordon of Old South Church in Boston; Washington Gladden, whose great success
came at the First Congregational Church of Columbus, Ohio; George Harris, professor of theology at
Andover and spokesman through \emph{The Andover Review}; and the great popularizer, Lyman Abbott, of
whom F.~H.~Foster said that probably no one in America did so much to move people toward a liberal
theology. Their allies included Henry Ward Beecher, Egbert Smyth, William Jewett Tucker and William Newton Clarke, among others. The movement was a protest, in the name of modern thought,
against the ``old theology'' of evangelical Puritanism, the New England Theology. Its last
outstanding spokesman, Edwards Amasa Park, resigned his chair at Andover in 1881, and Harris
succeeded him two years later, the watershed between old and new.

Before 1880 most New England ministers held the sovereignty of God, the innate depravity of mankind
traced to the first man's sin, the atonement of Christ as the ground of forgiveness, the Holy Spirit
as essential to conversion, and the eternal separation of saved and lost. After 1880 every one of
these came under fire. The most famous controversy broke out at Andover Seminary, founded by
Congregationalists in 1808 to counter the Unitarian drift of Harvard, whose founders had required
the faculty to subscribe to a Calvinist creed. By 1880 professors found this impossible and said so.
The spark was a series of articles in the \emph{Andover Review} in which Egbert C.~Smyth, Tucker and
Harris argued that heathen who die without knowing the gospel will be given the chance, after death and before the last judgment, to accept or refuse it. In the debate that followed the faculty moved
step by step to the open defense of liberal theology. The Board of Visitors charged Smyth with
departing from the creed, as a test case, and dismissed him; after years of maneuver the Supreme
Court of Massachusetts voided the dismissal in 1892. By then almost every American denomination had
its own celebrated heresy case.

\begin{keyfacts}
\keypeople{Lyman Abbott; Harry Emerson Fosdick; H.~Richard Niebuhr; Adolf von Harnack; Charles Lyell;
Charles Darwin; David Friedrich Strauss; Friedrich Schleiermacher; Albrecht Ritschl; Washington
Gladden; George Harris; Egbert C.~Smyth; Edwards Amasa Park.}
\keydates{%
  \kd{1830}{Lyell's \emph{Principles of Geology}}
  \kd{1834}{Death of Schleiermacher, ``father of modern theology''}
  \kd{1859}{Darwin's \emph{The Origin of Species}}
  \kd{1871}{Darwin's \emph{The Descent of Man}}
  \kd{1889}{Death of Albrecht Ritschl}
  \kd{1892}{Massachusetts court voids the dismissal of Egbert Smyth at Andover Seminary; Lyman Abbott's
    \emph{The Evolution of Christianity}}}
\keyterms{Protestant liberalism; evangelical and modernistic liberalism; immanence and transcendence;
evolution; natural selection; lower and higher criticism; historical Jesus; myth; Romanticism;
religious experience; God-consciousness; value judgments; kingdom of God on earth; New Theology.}
\end{keyfacts}
